\documentclass[10pt,twocolumn]{article}

\usepackage[margin=0.72in]{geometry}
\usepackage{amsmath}
\usepackage{booktabs}
\usepackage{graphicx}
\usepackage{microtype}
\usepackage{url}

\title{Auditing Normal-Reference Coordinate Drift in a Staged Skeleton JEPA}
\author{Anonymous authors}
\date{}

\begin{document}
\maketitle

\begin{abstract}
We audit a project-specific, S-JEPA-inspired skeleton model trained first on normal-labeled gait and then on four condition-labeled GAVD groups. The current experiment uses no added-normal dataset. After availability and pose-quality checks, it contains 626 sequences from 93 source videos. Across 270 matched normal sequences, the mean cosine between each later representation and its own Stage-0 representation falls from 0.700 after Stage 1 to 0.297 after Stage 4. Checkpoint reload reproduces the curve within $4.51\times10^{-7}$. This establishes raw coordinate drift for one seed, not catastrophic forgetting. A sequence-split classifier reaches 0.899 macro-F1, but every test row was encoder-exposed and 64 source videos cross the classifier split. We present the result as an auditable representation study and state the controls needed for functional-retention and generalization claims.
\end{abstract}

\section{Introduction}
Joint-Embedding Predictive Architectures learn by predicting target-encoder features rather than reconstructing every input value \cite{assran2023ijepa,bardes2024vjepa,abdelfattah2024sjepa}. In a normal-first model, later training can move the internal normal reference. A raw cosine audit can detect this movement, but it cannot by itself distinguish lost information from a shared rotation of the feature basis.

We make that distinction explicit for a gait skeleton JEPA variant. It is not a reproduction of the published S-JEPA: it uses a fixed 12-landmark target sampler, VICReg, and a later label-aware group term. Our verified contribution is a fingerprint-bound raw anchor curve that is recomputable from the saved checkpoints, plus a claim audit that separates coordinate drift, functional retention, in-corpus label readability, and unseen-source generalization.

\section{Method}
\subsection{Data and preprocessing}
The raw local GAVD inventory has 666 sequences from 103 source videos \cite{ranjan2025gavd}. Availability checks retain 642 sequences from 94 videos. Requiring at least 0.50 observed coverage over 12 gait-focused landmarks leaves the cohort in Table~\ref{tab:data}. These are dataset folder annotations, not diagnoses made by this project.

\begin{table}[t]
\centering
\small
\caption{Current GAVD-only cohort.}
\label{tab:data}
\begin{tabular}{lrr}
\toprule
Annotation & Sequences & Videos \\
\midrule
Normal & 270 & 29 \\
Parkinson's & 41 & 9 \\
Stroke & 74 & 18 \\
Myopathic & 183 & 28 \\
Cerebral palsy & 58 & 9 \\
\midrule
Total & 626 & 93 \\
\bottomrule
\end{tabular}
\end{table}

MediaPipe supplies 33 landmarks per frame \cite{grishchenko2022blazepose,mediapipePoseLandmarker}. The pipeline interpolates short internal gaps, centers on the pelvis, scales by body width, and resizes every sequence to 64 frames. Resizing removes native duration unless timing is supplied separately. Across the 642 availability-filtered pose caches, extraction labels are mixed, a potential provenance shortcut, although their recorded pose-model hash agrees.

\subsection{Model and curriculum}
A token contains one landmark over four frames. Only the shoulders, hips, knees, ankles, heels, and foot tips can be prediction targets. The model uses a view encoder, an exponential-moving-average target encoder, and a predictor. The loss is
\begin{equation}
\mathcal L=\mathcal L_{\mathrm{JEPA}}+0.05\mathcal L_{\mathrm{VICReg}}+0.25\mathcal L_{\mathrm{group}},
\end{equation}
where VICReg discourages collapse \cite{bardes2022vicreg}. The group term is off during normal-only Stage 0 and later uses folder labels to encourage compactness and centroid separation.

One model trains for 300 epochs on 270 normal rows. Four 75-epoch stages add Parkinson's, stroke, myopathic, and cerebral-palsy rows with balanced replay. Model, target, center, and projector state continue, while AdamW restarts at each stage. AdamW uses betas $(0.9,0.95)$, weight decay 0.05, learning rate $10^{-3}$ at Stage 0 and $3\times10^{-4}$ later, with four rows per active condition in each batch. The run has 600 epochs, 40,800 optimizer updates, seed 42, and the MPS backend. Exact hardware and deterministic settings were not recorded.

\subsection{Anchor and claim levels}
For the same 270 normal sequences at each stage, we report
\begin{equation}
a_t=\frac{1}{|N|}\sum_{x\in N}\frac{z_t(x)^\top z_0(x)}{\lVert z_t(x)\rVert_2\lVert z_0(x)\rVert_2}.
\end{equation}
This is the average of matched per-sequence cosines, not a cosine between cohort centroids. All anchor rows were used in training. The mean is sequence-weighted: the top two normal videos supply 105 of 270 rows and the top three supply 137. Population-level language therefore needs equal-video weighting and source-cluster uncertainty.

\section{Results}
The anchor values after the four later stages are 0.7002, 0.5021, 0.3962, and 0.2966. Reloading the checkpoints reproduces the saved summary to within $4.51\times10^{-7}$. Feature standard deviation remains 0.229 at Stage 4, arguing against total constant collapse.

\begin{figure}[t]
\centering
\includegraphics[width=\columnwidth]{figures/bbfm_drift_curve.pdf}
\caption{Raw matched-sequence anchor cosine. Stage 0 defines the reference. Alignment and functional controls are still required.}
\label{fig:drift}
\end{figure}

The final 384-dimensional training-corpus summaries have cosine silhouette 0.3617, minimum between-centroid distance 0.0863, and mean within-condition distance 0.0783. This structure is label-informed because the same folder labels shaped the later encoder objective.

\begin{table}[t]
\centering
\small
\caption{Descriptive five-class probes.}
\label{tab:readout}
\begin{tabular}{lrrr}
\toprule
Probe & Acc. & Bal. acc. & Macro-F1 \\
\midrule
All 626 rows & 0.920 & 0.900 & 0.899 \\
Exact 47/21 & 0.857 & 0.880 & 0.861 \\
Missingness only & 0.441 & 0.427 & 0.355 \\
\bottomrule
\end{tabular}
\end{table}

In the all-row probe, all 188 test rows were used by the encoder and 64 source videos cross the classifier split. In the exact probe, all 21 test rows were encoder-exposed and all nine test videos overlap classifier training. These scores are in-corpus readouts, not generalization estimates.

\section{Validity audit and required experiments}
We exclude the saved margin-ablation, AnchorGuard, grouped Lane C, and predictive-surprise result families. They load an old augmented checkpoint, lack complete current lineage, or combine old checkpoints with current rows. They cannot support mechanism, repair, or forecasting claims.

A strong drift or forgetting claim needs three to five full seeds, orthogonal-Procrustes alignment plus linear CKA or SVCCA, a same-update continued-normal control, held-out normal-function measurements at every checkpoint, and an order control. An unseen-source performance claim requires preprocessing, encoder training, and readout fitting inside every outer source-video fold \cite{roberts2017crossvalidation}.

\section{Discussion and conclusion}
The experiment establishes a substantial raw-coordinate change in matched normal sequences for one audited run: the mean cosine reaches 0.297. It does not establish that useful normal-gait information was lost. It also does not establish clinical or unseen-source performance. Source video is not a person identifier, folder labels are not independently adjudicated diagnoses, and camera, pose missingness, and extraction history may correlate with labels.

GAVD distributes annotations and public video URLs rather than raw files; independent retrieval must comply with YouTube terms, institutional ethics requirements, and applicable copyright, privacy, and data-protection rules \cite{gavdRepo2026}. This analysis uses derived poses and does not infer identity. A public artifact should not redistribute raw videos or identity-bearing frames. The institutional ethics determination and data-use review must be recorded before submission.

The methodological conclusion is direct: raw drift, functional retention, and downstream generalization are different endpoints. Reporting them separately prevents a checkpoint-recomputable telemetry signal from becoming an unsupported clinical or continual-learning claim.

\bibliographystyle{plain}
\bibliography{references}

\end{document}
